\section{Feasibility of Fairness Notions}
\label{sec:feasibility}

Feasible notions:

\begin{enumerate}
\item EF1 is polynomial-time feasible for equal entitlements when for each agent $i$,
    $v_i$ is doubly-monotone, i.e., the items $M$ can be partitioned into sets $G_i$ and $C_i$
    such that $v_i(g \mid S) \ge 0$ for all $S \subseteq M$ and $g \in G_i \setminus S$
    and $v_i(c \mid S) \le 0$ for all $S \subseteq M$ and $c \in C_i \setminus S$
    \cite[Theorem~4]{bhaskar2021approximate}.
\item EF1 is polynomial-time feasible for additive valuations over goods
    \cite[Theorem~3.3]{chakraborty2021weighted}.
\item EF1 is polynomial-time feasible for additive disutilities over chores
    \cite[Theorem~19]{springer2024almost}.
\item MMS is feasible for two agents and equal entitlements (using the cut-and-choose protocol).
\item WMMS is feasible for identical valuations.
\item PROPm is feasible for equal entitlements and additive valuations over goods
    \cite{baklanov2021propm}.
\item A PROP1+PO allocation always exists for additive valuations (over mixed manna)
    \cite{aziz2020polynomial}.
\item EFX is polynomial-time feasible for identical additive valuations over goods
    or over chores \cite[Theorem 5]{springer2024almost}.
\item EEFX is feasible for equal entitlements and cancelable valuations over goods
    or over chores \cite{caragiannis2022existence}.
\item EFX is feasible for equal entitlements and additive bivalued valuations over goods
    \cite[Theorem~4.1]{amanatidis2021maximum}.
\item MMS is feasible for equal entitlements and additive bivalued valuations over goods
    or over chores \cite{feige2022maximin}.
\item For equal entitlements and identical valuations, any leximin allocation is also GMMS
    (based on Theorem 3 of \cite{barman2018groupwise}).
\end{enumerate}

Infeasible notions:

\begin{enumerate}
\item PROP is infeasible for the instance in \cref{cex:single-item}.
\item PROP1 is infeasible for the instance in \cref{cex:ef-not-prop-submod}.
\item Remark 2 of \cite{babaioff2023fair} gives an example instance with two equally-entitled agents
    having identical submodular valuations with non-negative marginals for which no APS allocation exists.
    The instance can be converted to additionally have positive marginals.
\item Lemma C.1 of \cite{babaioff2023fair} gives an example instance with 3 equally-entitled agents
    having identical additive valuations with positive marginals for which no APS allocation exists.
\item An example with 3 equally-entitled agents having identical valuations with positive marginals
    for which no MMS allocation exists \cite{feige2022tight}.
    Same result also holds for negative marginals instead.
\item Consider an instance with 3 equally-entitled agents having identical additive valuations over 2 goods
    with values 10 and 1, respectively. Then no PROPx allocation exists.
\item For two agents having additive valuations over goods or over chores, an EFX allocation
    doesn't exist (Theorems 8 and 17 of \cite{springer2024almost}).
\end{enumerate}
